\documentclass[11pt, aspectratio=169]{beamer}

% --- PACKAGES ---
\usepackage[utf8]{inputenc}
\usepackage[french]{babel}
\usepackage{amsmath}
\usepackage{amssymb}
\usepackage{graphicx}
\usepackage{lmodern}
\usepackage{hyperref}


% --- THEME ---
\usetheme{Madrid}
\usecolortheme{default}
\definecolor{UVSQBlue}{RGB}{0, 51, 102}
\setbeamercolor{structure}{fg=UVSQBlue}
\setbeamercolor{frametitle}{bg=UVSQBlue, fg=white}
\setbeamercolor{title}{bg=UVSQBlue, fg=white}
\setbeamertemplate{navigation symbols}{}

% --- FOOTER ---
\setbeamertemplate{footline}{
  \leavevmode%
  \hbox{%
    \begin{beamercolorbox}[wd=.33\paperwidth,ht=2.25ex,dp=1ex,left]{author in head/foot}%
      \hspace{1em}\tiny Master 1 CHPS%
    \end{beamercolorbox}%
    \begin{beamercolorbox}[wd=.34\paperwidth,ht=2.25ex,dp=1ex,center]{title in head/foot}%
      \tiny Simulation d'un Banc de Poissons%
    \end{beamercolorbox}%
    \begin{beamercolorbox}[wd=.33\paperwidth,ht=2.25ex,dp=1ex,right]{date in head/foot}%
      \tiny Décembre 2025\hspace{1em}%
    \end{beamercolorbox}%
  }%
  \vskip0pt%
}

% --- INFOS ---
\title{Modélisation et Simulation distribuée du mouvement d'un Banc de Poissons}
\subtitle{Modélisation Multi-Agents avec les Règles de Reynolds}
\author{Yohann FRONT-REIGNIER \and Iram MADANI-FOUATIH \and Laurence HUANG}
\institute{Master 1 CHPS -- UVSQ Paris-Saclay}
%\date{Décembre 2025}
\titlegraphic{
\includegraphics[height=2cm]{ISTY-logo-2025.png}}

\begin{document}

\begin{frame}
\titlepage
\end{frame}

% --- SLIDE 2 ---
\begin{frame}
\frametitle{Contexte}
\begin{itemize}
    \setlength\itemsep{1em}
    \item Simuler les comportements collectifs émergents
    \item Bancs de poissons, flocks d'oiseaux
    \item 3 règles locales simples appliquées à chaque agent
    \item \textbf{Objectifs S1} : implémentation séquentielle + visualisation 2D
    \item \textbf{Semestre 2} : parallélisation OpenMP/MPI
\end{itemize}
\end{frame}

% --- SLIDE 3 ---
\begin{frame}
\frametitle{Les 3 Règles de Reynolds}
\begin{enumerate}
    \setlength\itemsep{1em}
    \item \textbf{Séparation} : Éviter les collisions avec les voisins proches
    \item \textbf{Alignement} : S'adapter à la vélocité moyenne du groupe local
    \item \textbf{Cohésion} : Se diriger vers le centre de masse des voisins
\end{enumerate}
\vspace{1cm}
\centering
\textit{Appliquées localement $\Rightarrow$ comportement global complexe et réaliste}
\end{frame}

% --- SLIDE 4 ---
\begin{frame}
\frametitle{Formulation Mathématique}
Pour chaque boid $b$, on calcule trois forces :
\vspace{0.3cm}

\textbf{Séparation :}
$$\vec{F}_{sep} = -\sum_{i \in N(b)} \frac{\vec{p}_b - \vec{p}_i}{\|\vec{p}_b - \vec{p}_i\|^2}$$

\textbf{Alignement :}
$$\vec{F}_{ali} = \frac{1}{|N(b)|}\sum_{i \in N(b)} \vec{v}_i - \vec{v}_b$$

\textbf{Cohésion :}
$$\vec{F}_{coh} = \frac{1}{|N(b)|}\sum_{i \in N(b)} \vec{p}_i - \vec{p}_b$$

\vspace{0.2cm}
\small
Mise à jour : $\vec{a} = \vec{F}_{total}/m$, $\vec{v} \leftarrow \vec{v} + \vec{a}\Delta t$, $\vec{p} \leftarrow \vec{p} + \vec{v}\Delta t$
\end{frame}

% --- SLIDE 5 ---
\begin{frame}
\frametitle{Architecture Logicielle}
\begin{itemize}
    \setlength\itemsep{0.8em}
    \item \textbf{Vector2D} : Opérations vectorielles (position, vitesse, accélération)
    \item \textbf{Boid} : Agent autonome avec les 3 règles Reynolds
    \item \textbf{Flock} : Gestionnaire de collection de boids
    \item \textbf{SpatialGrid} : Grille d'accélération pour recherche de voisins
    \item \textbf{GUI (SFML)} : Interface interactive avec sliders temps réel
\end{itemize}
\end{frame}

% --- SLIDE 6  ---
\begin{frame}
\frametitle{Optimisation : SpatialGrid}

\vspace{0.5cm}

\textbf{\large 1. Le Problème} \\
Recherche de voisins naïve = $O(n^2)$ par frame. \textbf{Trop lent.}

\vspace{1cm}

\textbf{\large 2. La Solution : Partitionnement}
\begin{itemize}
    \setlength\itemsep{0.5em}
    \item Partitionner l'espace en cellules = $O(k)$ (où $k$ = boids/cellule)
    \item \textbf{Implémentation} : Grille 2D de vecteurs \texttt{Boid*}
    \item \textbf{Gain} : 500 boids $\rightarrow$ 5x plus rapide (50k comparaisons vs 250k)
\end{itemize}

\end{frame}
% --- SLIDE 7 ---
\begin{frame}
\frametitle{Benchmarks de performance}

\begin{itemize}
\setlength\itemsep{0.8em}

\item \textbf{Temps vs rayon (N = 500)} :
le SpatialGrid reste nettement plus rapide que l'algorithme naïf,
et l'écart se creuse lorsque le rayon de voisinage augmente.

\item \textbf{Temps vs nombre de boids (rayon = 50 px)} :
la courbe naïve suit une croissance quadratique en $N$,
tandis que le SpatialGrid se rapproche d'une croissance quasi linéaire.

\item Benchmarks obtenus en moyennant le temps par itération
sur plusieurs centaines de frames, afin de lisser les fluctuations
et comparer clairement les deux approches.
\end{itemize}

%\vspace{0.5cm}
%\centering
\includegraphics[width=0.42\textwidth]{/Users/nouzhafou/Mod-lisation-et-simulation-distribu-e-du-mouvement-d-un-banc-de-poissons/figures/bench_time_vs_radius.png}\hfill
\includegraphics[width=0.42\textwidth]{/Users/nouzhafou/Mod-lisation-et-simulation-distribu-e-du-mouvement-d-un-banc-de-poissons/figures/bench_time_vs_N.png}

\end{frame}
% --- SLIDE 7 ---
\begin{frame}
\frametitle{Détails d'Implémentation}
\begin{itemize}
    \setlength\itemsep{0.8em}
    \item \textbf{Langage} : C++20 (moderne, performant)
    \item \textbf{Build} : CMake 3.20+, GCC/Clang avec flags \texttt{-O3}
    \item \textbf{Graphisme} : SFML 2.6.1 pour visualisation 2D
    \item \textbf{Tests} : Google Test, CTest automatique
    \item \textbf{Documentation} : Doxygen
    \item \textbf{Git} : Gestion de version (97 commits)
\end{itemize}
\end{frame}


% --- SLIDE 8 ---
\begin{frame}
\frametitle{Résultats Observés}
\begin{itemize}
    \setlength\itemsep{0.8em}
    \item \textbf{Cohésion} : Agents restent groupés, structure maintenue
    \item \textbf{Alignement} : Direction moyenne émerge rapidement
    \item \textbf{Séparation} : Collisions évitées efficacement
    \item \textbf{Patterns} : Tourbillons, colonnes, vagues observés
    \item \textbf{Performance} : \textbf{60 FPS stable} avec 500 boids
\end{itemize}
\end{frame}

% --- SLIDE 9 ---
\begin{frame}
\frametitle{Interface Interactive (SFML)}
\begin{itemize}
    \setlength\itemsep{0.8em}
    \item Sliders temps réel : Séparation, Alignement, Cohésion (0-10)
    \item Rayon de voisinage : 5-200 pixels
    \item Nombre de boids : 10-500
    \item Bouton \textit{Apply} pour réinitialiser la simulation
    \item Affichage des FPS
\end{itemize}
\end{frame}

% --- SLIDE 10 ---
\begin{frame}
\frametitle{Contributions de l'équipe}
\begin{columns}[T]
    \begin{column}{0.33\textwidth}
        \textbf{Yohann} \\ (36 commits)
        \footnotesize
        \begin{itemize}
            \item Vector2D
            \item Architecture
            \item Rapport
            \item Doc Doxygen
        \end{itemize}
    \end{column}
    \begin{column}{0.33\textwidth}
        \textbf{Iram} \\ (38 commits)
        \footnotesize
        \begin{itemize}
            \item Boid (Règles)
            \item GUI SFML
            \item Tests Boid
        \end{itemize}
    \end{column}
    \begin{column}{0.33\textwidth}
        \textbf{Laurence} \\ (8 commits)
        \footnotesize
        \begin{itemize}
            \item Flock
            \item SpatialGrid
            \item Optimisations
            \item Benchmarks
        \end{itemize}
    \end{column}
\end{columns}
\end{frame}

% --- SLIDE 11 ---
\begin{frame}
\frametitle{Tests et Validation}
\textbf{Tests Unitaires (Google Test) :}
\begin{itemize}
    \item \texttt{test\_vector.bin} : Opérations, magnitude, normalisation
    \item \texttt{test\_boid.bin} : Forces, mise à jour, règles Reynolds
    \item \texttt{test\_flock.bin} : Gestion collection, SpatialGrid
    \item \texttt{test\_ui.bin} : Intégration GUI
\end{itemize}
\vspace{0.5cm}
\textbf{Validation :}
\begin{itemize}
    \item Validation visuelle des 3 règles émergentes
    \item Tous les tests passent (100\% succès)
\end{itemize}
\end{frame}

% --- SLIDE 12 ---
\begin{frame}
\frametitle{Perspectives (Semestre 2)}
\begin{itemize}
    \setlength\itemsep{1em}
    \item \textbf{Parallélisation} : OpenMP (threads), MPI (cluster)
    \item \textbf{Extensions} : 3D (OpenGL), obstacles, prédateurs/proies
    \item \textbf{Optimisations} : kd-tree, SIMD, calcul sur GPU
    \item \textbf{Benchmarks} : Scalabilité, performance comparative
\end{itemize}
\vspace{0.5cm}
\textbf{État actuel :} Code propre, testé, prêt pour l'optimisation.
\end{frame}

\begin{frame}
\frametitle{Conclusion}

\begin{itemize}
\setlength\itemsep{1em}
\item Simulateur complet et fonctionnel (3 règles de Reynolds)
\item Architecture modulaire et optimisée (SpatialGrid)
\item Visualisation fluide (60 FPS @ 500 agents)
\item Compétences acquises : C++ moderne, CMake, Git, Systèmes Multi-Agents
\item Code source disponible : \href{https://github.com/iramyas/Mod-lisation-et-simulation-distribu-e-du-mouvement-d-un-banc-de-poissons.git}{\textcolor{blue}{lien vers GitHub}}
\end{itemize}

\end{frame}


\begin{frame}
\centering
\Large \textcolor{UVSQBlue}{\textbf{Merci de votre attention !}} \\
\vspace{0.5cm}
\end{frame}

\end{document}
